\section*{CHƯƠNG 2: THỰC TRẠNG PHÂN HÓA LẬP TRƯỜNG VÀ NHỮNG CHUYỂN BIẾN TƯ TƯỞNG CỦA TẦNG LỚP NHO SĨ THẾ KỶ XVIII}
\addcontentsline{toc}{section}{CHƯƠNG 2: THỰC TRẠNG PHÂN HÓA LẬP TRƯỜNG VÀ NHỮNG CHUYỂN BIẾN TƯ TƯỞNG CỦA TẦNG LỚP NHO SĨ THẾ KỶ XVIII}

\subsection*{2.1. Sự phân hóa thái độ chính trị -- xã hội của tầng lớp nho sĩ trước các biến cố lịch sử}
\addcontentsline{toc}{subsection}{2.1. Sự phân hóa thái độ chính trị -- xã hội của tầng lớp nho sĩ trước các biến cố lịch sử}

\subsubsection*{2.1.1. Xu hướng bảo thủ giữ chữ trung quân mù quáng và xu hướng bất mãn ẩn dật thúc đẩy phong trào tư học}
\addcontentsline{toc}{subsubsection}{2.1.1. Xu hướng bảo thủ giữ chữ trung quân mù quáng và xu hướng bất mãn ẩn dật thúc đẩy phong trào tư học}

Trước sự cáo chung của chính quyền Lê -- Trịnh và chúa Nguyễn, mâu thuẫn giữa tín điều ``trung quân'' với một dòng họ và nghĩa vụ đối với vận mệnh dân tộc đã chia rẽ tầng lớp nho sĩ thành những thái độ ứng xử trái ngược nhau.

Ở cực bảo thủ nhất là bộ phận nho sĩ chịu sự chi phối nặng nề của giáo điều ``Trung thần bất sự nhị quân'' (Bề tôi trung không thờ hai chúa) và thuyết ``Chính danh'' Tống Nho. Đối với họ, ngôi vị nhà Lê là thiên mệnh bất biến, còn mọi cuộc nổi dậy của nông dân đều bị quy thành hành vi tiếm nghịch. Khi trật tự cũ sụp đổ trước sức tiến công của quân Tây Sơn, nhiều người rơi vào bế tắc tư tưởng và chọn cách giữ tiết cực đoan: người thì tự vẫn tuẫn tiết theo triều cũ như Lý Trần Quán (tự chôn sống năm 1786), Trần Công Xán (chịu chết khi đi sứ Phú Xuân năm 1787); người thì bôn tẩu cần vương bất thành rồi lui về ẩn cư, cự tuyệt hợp tác với triều đại mới đến trọn đời như Trần Danh Án, Lê Duy Đản.

Nghiêm trọng hơn, sự mù quáng về ý thức hệ đã đẩy nhóm quan lại nho sĩ như Lê Quynh bước qua ranh giới phản bội dân tộc. Nhân danh khôi phục cơ nghiệp nhà Lê, nhóm này đã hộ tống vua Lê Chiêu Thống sang cầu viện triều đình Mãn Thanh, trực tiếp dẫn đường cho quân xâm lược của Tôn Sĩ Nghị tràn vào giày xéo đất nước. Dưới góc độ duy vật lịch sử, đây là bi kịch tất yếu khi lòng trung thành giai cấp hẹp hòi bị đặt cao hơn chủ quyền quốc gia và quyền sống của nhân dân.

Khác với xu hướng bảo thủ trên, một bộ phận đông đảo nho sĩ có lương tri lại chọn con đường thoái lui khỏi chốn quan trường. Họ chán ghét sự mục nát, phe phái trong phủ chúa cung vua, song do giới hạn giai cấp nên vẫn dè dặt, hoài nghi trước phong trào nông dân Tây Sơn. Vận dụng triết lý ``Độc thiện kỳ thân'' (giữ mình trong sạch giữa thời loạn) và ``Dĩ thân tàng'' (ẩn mình chờ thời), nhiều danh nho từ quan trở về quê nhà như Bùi Huy Bích, Phạm Quý Thích, hay La Sơn Phu Tử Nguyễn Thiếp trong giai đoạn ẩn cư tại trại Bùi Phong trên dãy núi Thiên Nhẫn.

Sự thoái lui chính trị của các bậc túc nho này lại mang đến một hệ quả văn hóa tích cực là thúc đẩy mạng lưới tư học phát triển mạnh mẽ. Giữa lúc Quốc Tử Giám và trường học công lập suy tàn, các lớp học tư gia của Bùi Huy Bích ở Định Công hay của Nguyễn Thiếp ở xứ Nghệ đã thu hút hàng trăm môn sinh đến thụ giáo. Tại đây, người thầy không chỉ truyền dạy kinh sử mà còn đề cao đạo làm người, giữ gìn sĩ khí trước tệ nạn mua danh bán tước nơi trường thi. Phong trào tư học thế kỷ XVIII nhờ đó đã bảo toàn mạch nguồn văn hiến và lan tỏa tri thức xuống tận cấp cơ sở làng xã.

\subsubsection*{2.1.2. Xu hướng tiến bộ dấn thân cùng phong trào Tây Sơn và nét đặc thù khai phóng của nho sĩ Đàng Trong}
\addcontentsline{toc}{subsubsection}{2.1.2. Xu hướng tiến bộ dấn thân cùng phong trào Tây Sơn và nét đặc thù khai phóng của nho sĩ Đàng Trong}

Vượt lên trên tâm lý hoài cổ hay ẩn dật, bộ phận nho sĩ ưu tú nhất của Bắc Hà đã dứt khoát rũ bỏ định kiến giai cấp để đồng hành cùng phong trào Tây Sơn. Nhận thức của họ bắt nguồn từ sự kế thừa tư tưởng ``Dân vi bản'' (Dân là gốc của nước): triều đại nào bảo vệ được độc lập dân tộc và đem lại cuộc sống yên ổn cho trăm họ thì triều đại đó nắm giữ tính chính thống lịch sử. Đáp lại tấm lòng trân trọng hiền tài của Hoàng đế Quang Trung qua bản \textit{Chiếu cầu hiền} (do Ngô Thì Nhậm chấp bút) và tư tưởng cải cách trong \textit{Chiếu lập học} -- ``Xây dựng đất nước lấy việc dạy học làm đầu, tìm lẽ trị bình lấy việc tuyển nhân tài làm gấp'' -- nhiều sĩ phu tài danh đã đem trí tuệ phụng sự vương triều mới.

Tiêu biểu nhất cho xu hướng nhập thế này là Ngô Thì Nhậm (1746--1803). Mùa đông năm 1788, trước thế giặc Mãn Thanh ồ ạt tiến công, chính ông đã đề xuất kế sách rút quân chiến lược khỏi Thăng Long về lập phòng tuyến liên hoàn Tam Điệp -- Biện Sơn nhằm bảo toàn chủ lực, làm kiêu lòng giặc, tạo bàn đạp cho cuộc hành quân thần tốc đại phá quân Thanh mùa Xuân năm Kỷ Dậu 1789. Sau ngày chiến thắng, được vua Quang Trung giao trọng trách bang giao, Ngô Thì Nhậm đã vận dụng sách lược ngoại giao vừa kiên quyết giữ vững chủ quyền, vừa mềm dẻo khéo léo, buộc triều đình Càn Long phải bãi bỏ ý định phục thù và chính thức công nhận vương triều Tây Sơn.

Sát cánh cùng Ngô Thì Nhậm là những bậc danh nho đức độ và thực tài. La Sơn Phu Tử Nguyễn Thiếp (1723--1804), người từng nhiều lần khước từ lời vời của chúa Trịnh, đã cảm phục tầm vóc anh hùng và sự cầu thị của vua Quang Trung mà nhận chức Viện trưởng Viện Sùng Chính, chủ trì dịch kinh sách chữ Hán sang chữ Nôm để cải cách giáo dục và tham mưu quy hoạch Phượng Hoàng Trung Đô tại Nghệ An. Trong khi đó, Phan Huy Ích (1751--1822) đảm nhận soạn thảo văn thư bang giao, xây dựng điển chế triều nghi và để lại bộ \textit{Dụ Am văn tập} có giá trị sử liệu cao.

Ở phía Nam, tầng lớp nho sĩ xứ Đàng Trong (đặc biệt tại vùng Gia Định -- Nam Bộ) lại định hình phong cách khai phóng rõ nét. Do sinh trưởng trên vùng đất mới khai phá, nơi hội tụ cư dân đa nguồn gốc cùng hoạt động thương cảng sôi động, kẻ sĩ phương Nam ít bị ràng buộc bởi lễ giáo từ chương khắt khe như Bắc Hà, mà giàu tính thực tiễn, trọng nghĩa khí và gần gũi với dân nghèo.

Gương mặt đại diện cho sĩ phu Nam Bộ thế kỷ XVIII là danh sư Võ Trường Toản (mất năm 1792), bậc ``Sùng Đức xử sĩ'' được người đương thời tôn vinh là ``Sư biểu đất Gia Định''. Giữa cuộc giao tranh quyết liệt giữa Nguyễn Ánh và Tây Sơn, ông giữ thái độ độc lập, không ra làm quan mà mở trường dạy học tại làng Hòa Hưng. Võ Trường Toản đề xướng phương châm giáo dục ``Tri ngôn dưỡng khí'' (hiểu thấu đạo lý để nuôi dưỡng chí khí khẳng khái), chú trọng rèn luyện thực đức hơn lối văn chương cử nghiệp. Trường học của ông đã đào tạo nên lớp trí thức trụ cột của miền Nam, nổi bật là nhóm ``Gia Định tam gia'' gồm Trịnh Hoài Đức, Lê Quang Định và Ngô Nhân Tịnh -- những tác giả đã biên soạn bộ địa chí \textit{Gia Định thành thông chí} và góp phần hun đúc nên ``Hào khí Đồng Nai''.

\subsection*{2.2. Những bước chuyển biến trong thế giới quan triết học, khuynh hướng học thuật và ý thức nhân văn}
\addcontentsline{toc}{subsection}{2.2. Những bước chuyển biến trong thế giới quan triết học, khuynh hướng học thuật và ý thức nhân văn}

\subsubsection*{2.2.1. Bước chuyển từ Lý học Tống Nho kinh viện sang trào lưu thực học và tư duy bách khoa khảo chứng}
\addcontentsline{toc}{subsubsection}{2.2.1. Bước chuyển từ Lý học Tống Nho kinh viện sang trào lưu thực học và tư duy bách khoa khảo chứng}

Suốt từ thế kỷ XV, dòng Tống Nho (phái Trình -- Chu) chi phối học thuật Đại Việt với xu hướng thiên về bình giải kinh viện các phạm trù trừu tượng như Tâm, Tính, Mệnh và quy phạm hóa đạo đức tam cương ngũ thường. Bước sang thế kỷ XVIII, lối học tư biện sáo rỗng ấy tỏ ra bất lực trước hàng loạt vấn đề cấp bách của đất nước như trị thủy, cứu đói, quản lý điền chế và phát triển giao thương. Từ yêu cầu thực tiễn đó, giới nho sĩ tiến bộ đã chuyển hướng mạnh mẽ sang trào lưu thực học và phương pháp khảo chứng bách khoa, đề cao tri thức gắn với ``quốc kế dân sinh''.

Đỉnh cao của khuynh hướng thực học và tư duy triết học thế kỷ XVIII kết tinh ở nhà bác học Lê Quý Đôn (1726--1784) trên bốn phương diện chủ yếu:

Về bản thể luận, trong mục \textit{Lý khí} của bộ sách \textit{Vân đài loại ngữ}, Lê Quý Đôn dùng phạm trù ``Khí'' để cắt nghĩa cội nguồn vật chất của vũ trụ. Ông quan niệm Khí là thực thể nguyên sơ lấp đầy không gian, sự tụ tán và vận động tự thân của âm dương nhị khí tạo thành muôn vật trong trời đất. Luận điểm này mang khuynh hướng duy vật tự phát rõ rệt, đối lập với cách giải thích thế giới bằng ý chí thần linh siêu nhiên.

Về quan hệ giữa Lý và Khí, Lê Quý Đôn trực tiếp phản bác mệnh đề duy tâm khách quan của Chu Hy (vốn đặt Lý có trước và sinh ra Khí). Ông khẳng định ``Lý ở trong Khí'', coi Lý không phải là thực thể tinh thần tồn tại biệt lập bên ngoài vật chất, mà chỉ là quy luật xác định vốn có của từng sự vật hữu hình. Do quy luật nằm ngay trong sự vật, muốn nhận thức đúng đắn thì con người phải khảo sát trực tiếp bản thân thế giới hiện thực.

Về phương pháp nhận thức, Lê Quý Đôn giải thích lại mệnh đề ``cách vật trí tri'' theo hướng thực chứng: nghiên cứu phải dựa trên thu thập chứng cứ, đối chiếu thư tịch cổ kim và khảo sát thực địa. Ông để lại hệ thống công trình bách khoa đồ sộ gồm \textit{Vân đài loại ngữ}, \textit{Kiến văn tiểu lục}, \textit{Đại Việt thông sử} và \textit{Phủ biên tạp lục}. Riêng trong \textit{Phủ biên tạp lục}, những trang ghi chép tỉ mỉ của ông về vị trí địa lý, sản vật cùng hoạt động thực thi chủ quyền của đội Hoàng Sa và đội Bắc Hải đã trở thành nguồn sử liệu gốc xác thực về chủ quyền của Việt Nam đối với hai quần đảo Hoàng Sa, Trường Sa.

Về tư duy kinh tế, Lê Quý Đôn vượt qua định kiến ``trọng nông ức thương'' truyền thống. Qua các khảo cứu trong mục \textit{Điển vựng}, \textit{Phẩm vật} sách \textit{Vân đài loại ngữ} và \textit{Phủ biên tạp lục}, ông chủ trương kết hợp hài hòa giữa phát triển nông nghiệp với mở mang thủ công nghiệp, khai mỏ, lưu thông tiền tệ -- thương mại và bồi dưỡng tầng lớp hiền tài. Tinh thần kinh tế toàn diện đó thường được hậu thế đúc kết qua câu phương ngôn ``Phi nông bất ổn, phi công bất phú, phi thương bất hoạt, phi trí bất hưng'', phản ánh nhận thức biện chứng về bốn trụ cột giữ cho quốc gia ổn định, giàu mạnh và hưng thịnh.

\subsubsection*{2.2.2. Khuynh hướng hòa đồng Tam giáo và sự trỗi dậy của chủ nghĩa nhân đạo trong sáng tác của nho sĩ}
\addcontentsline{toc}{subsubsection}{2.2.2. Khuynh hướng hòa đồng Tam giáo và sự trỗi dậy của chủ nghĩa nhân đạo trong sáng tác của nho sĩ}

Song hành cùng trào lưu thực học là sự phục hưng tinh thần dung hợp Tam giáo (Nho -- Phật -- Đạo). Khi khuôn mẫu giáo điều Nho gia không còn giải đáp thỏa đáng những bi kịch nhân sinh thời loạn, nhiều nho sĩ đã tìm đến triết lý từ bi, vô ngã của Phật giáo và quan niệm thuận theo tự nhiên của Lão -- Trang để cân bằng đời sống tinh thần.

Công trình tiêu biểu nhất của xu hướng này là chuyên luận triết học \textit{Trúc Lâm tông chỉ nguyên thanh} do Ngô Thì Nhậm biên soạn dưới pháp danh Hải Lượng đại thiền sư. Kế thừa tư tưởng nhập thế ``Cư trần lạc đạo'' của Thiền phái Trúc Lâm thời Trần, Ngô Thì Nhậm chủ trương thống nhất Nho với Phật trên cùng một bản thể tâm tính: người trí thức khi xuất thế gánh vác việc quân quốc thì hành xử theo đạo Nho, khi quay về tu dưỡng nội tâm để giữ sự thanh tịnh, sáng suốt thì thấm nhuần đạo Thiền. Cách tiếp cận đó đã phá vỡ ranh giới bài xích tôn giáo hẹp hòi của Tống Nho từng tồn tại suốt nhiều thế kỷ.

Sự cởi mở về thế giới quan đồng thời khơi nguồn cho trào lưu chủ nghĩa nhân đạo phát triển rực rỡ trong văn học chữ Nôm của giới nho sĩ. Văn chương thế kỷ XVIII thoát dần khỏi công thức giáo huấn ``văn dĩ tải đạo'' để hướng thẳng vào số phận con người trần thế, lên án chiến tranh phong kiến và đòi quyền sống, quyền hạnh phúc cá nhân.

Mỗi tác phẩm lớn thời kỳ này đều mang giá trị phê phán hiện thực và cảm quan nhân văn sắc nét. Khúc ngâm \textit{Chinh phụ ngâm khúc} (nguyên tác chữ Hán của Đặng Trần Côn, bản diễn Nôm truyền tụng gắn với Đoàn Thị Điểm và Phan Huy Ích) là tiếng nói xót thương tuổi xuân người phụ nữ bị chia lìa bởi chiến tranh phi nghĩa. Trong khi đó, \textit{Cung oán ngâm khúc} của Nguyễn Gia Thiều phơi bày bi kịch tù túng nơi cung cấm để phê phán lối sống hưởng lạc của bậc đế vương. Ở mạch thơ Nôm truyền tụng, Hồ Xuân Hương dùng tiếng cười trào phúng đả kích trực diện thói đạo đức giả của giới nho sĩ hủ lậu và khẳng định khát vọng bình đẳng của người phụ nữ; còn truyện thơ \textit{Sơ kính tân trang} của Phạm Thái lại cất lên tiếng nói bảo vệ tình yêu tự do trước lễ giáo môn đăng hộ đối khắc nghiệt. Những thành tựu ấy đã tạo tiền đề nghệ thuật và tư tưởng trực tiếp cho sự ra đời kiệt tác \textit{Đoạn trường tân thanh} (Truyện Kiều) của Nguyễn Du vào đầu thế kỷ XIX.

Như vậy, biến động xã hội thế kỷ XVIII đã thúc đẩy tầng lớp nho sĩ Việt Nam thực hiện một cuộc tự vượt bỏ trên cả hai bình diện chính trị và học thuật. Từ sự phân hóa thái độ xuất xử đến sự hình thành trào lưu thực học, dung hợp Tam giáo và chủ nghĩa nhân đạo, diện mạo tinh thần của kẻ sĩ Đại Việt đã đạt tới độ chín muồi mới trong lịch sử tư tưởng dân tộc.
